% Wave-16 V2/20 --- Monstrous moonshine $j(\tau)$ and $V^\natural$
% via the two-stage factorisation $\Phi_3 = \SpCh_{\Sigma_2,C} \circ \PhiFA_3$
% Voices: Frenkel (VOA) + Carnahan (generalised moonshine)
% Vol III, Calabi--Yau quantum groups (Raeez Lorgat)

\documentclass[11pt]{article}
\usepackage[localtheorems]{raeez-math-template}

\providecommand{\Z}{\mathbb{Z}}
\providecommand{\Q}{\mathbb{Q}}
\providecommand{\C}{\mathbb{C}}
\providecommand{\HH}{\mathbb{H}}
\providecommand{\Lat}{\Lambda}
\providecommand{\Leech}{\Lambda_{24}}
\providecommand{\IIlor}{\mathrm{II}_{25,1}}
\providecommand{\IIonel}{\mathrm{II}_{1,1}}
\providecommand{\Vnat}{V^{\natural}}
\providecommand{\VLam}{V_{\Leech}}
\providecommand{\Mon}{\mathbb{M}}
\providecommand{\mMon}{\mathfrak{m}}
\providecommand{\Aut}{\mathrm{Aut}}
\providecommand{\Fcal}{\mathcal{F}}
\providecommand{\Sp}{\mathrm{Sp}}
\providecommand{\SpCh}{\mathrm{Sp}^{\mathrm{ch}}}
\providecommand{\PhiFA}{\Phi^{\mathrm{FA}}}
\providecommand{\EdHolFA}{E_d\text{-}\mathrm{hFA}}
\providecommand{\ETHolFA}{E_3\text{-}\mathrm{hFA}}
\providecommand{\kch}{\kappa_{\mathrm{ch}}}
\providecommand{\kcat}{\kappa_{\mathrm{cat}}}
\providecommand{\kBKM}{\kappa_{\mathrm{BKM}}}
\providecommand{\Tg}{T_g}
\providecommand{\Tgh}{T_{(g,h)}}
\providecommand{\ClaimStatusTheorem}{[Theorem]}
\providecommand{\ClaimStatusConjectured}{[Conjectured]}

\theoremstyle{plain}
\newtheorem{theorem}{Theorem}
\newtheorem{proposition}[theorem]{Proposition}
\newtheorem{lemma}[theorem]{Lemma}
\theoremstyle{definition}
\newtheorem{construction}[theorem]{Construction}
\newtheorem{conjecture}[theorem]{Conjecture}

\title{Moonshine, $\Vnat$, and $j(\tau)$ as witnesses of the
two-stage factorisation}
\author{Wave-16 V2/20\\ Frenkel + Carnahan voices}
\date{}

\begin{document}
\maketitle

\paragraph{Orientation.} The sibling note S3 locates the Borcherds
denominator $j(p)-j(q)$ at the $(T^{24}_\Leech,\mathrm{pt})$ cusp of the
two-stage factorisation $\Phi_3=\SpCh_{\Sigma_2,C}\circ\PhiFA_3$. V2
reads the same diagram from the vertex-algebra side: the graded
character $J(\tau)=j(\tau)-744$ and its twists $\Tg(\tau)$,
$\Tgh(\tau)$ are specialisation witnesses of $\PhiFA_3(X_B)$ at a
tower of $(\Sigma_2,C)$-loci.

\section*{1. $\Vnat$ as $\SpCh_{\Sigma_2,C}$ of the FLM $E_3$-hFA}

The Frenkel--Lepowsky--Meurman $\Z/2$-orbifold
\(
\Vnat=(\VLam)^{\Z/2}\oplus(\VLam^T)^{\Z/2}
\)
has central charge $c=24$, partition function $J(\tau)$, and
$\Aut\Vnat=\Mon$ (FLM 1988 Thm.~12.3.1; Griess 1982 for the
algebra). On the CY$_3$ $X_B=(T^{24}_\Leech\times E_\tau)/(\Z/2)$
the first-stage functor produces
\[
\PhiFA_3(X_B)\in\ETHolFA(X_B),
\]
an $E_3$-holomorphic factorisation algebra pinned up to
contractible choice by the Gerstenhaber bracket of degree $-2$ on
$\HH^\bullet(D^b X_B)$ and Costello--Gwilliam--Li locality.

\begin{construction}[$\Vnat$ as second-stage specialisation]
\label{cons:vnat-two-stage}
Take $\Sigma_2=T^{24}_\Leech/(\Z/2)$ (the Leech fibre, with the
FLM involution) and $C=E_\tau\subset X_B$. Then
\[
\SpCh_{\Sigma_2,E_\tau}\bigl(\PhiFA_3(X_B)\bigr)
\;=\;\Vnat \quad\text{as $E_1$-chiral algebra on $E_\tau$}.
\]
Factorisation homology $\int_{\Sigma_2}\PhiFA_3(X_B)$ is the
Dijkgraaf--Witten orbifold holomorphic-factorisation integral of
the Leech sigma model; the untwisted summand gives $(\VLam)^{\Z/2}$
and the $\mathrm{Aut}(\Sigma_2)^*(\Z/2)$-twist sector gives the
FLM twisted module $(\VLam^T)^{\Z/2}$. Restriction to $E_\tau$
returns the holomorphic VOA whose graded trace is $J(\tau)$.
\end{construction}

\noindent This is the precise Vol III sense in which the Monster
module is a chiral-algebra shadow of a CY$_3$: $\Vnat$ lives at
stage-2, while the $E_3$-hFA $\PhiFA_3(X_B)$ of stage-1 carries the
full moonshine datum. The native operadic level at $d=3$ is
$n_{\mathrm{native}}(3)=1$ (Prop.\ native-$E_n$-level,
\S e1-chiral): the $E_2$-braided structure of $\Vnat$ sits on
$\mathcal{Z}(\mathrm{Rep}^{E_1}(\Vnat))$, not on $\Vnat$ itself.

\section*{2. McKay--Thompson series as $g$-twisted specialisations}

For $g\in\Mon$ the McKay--Thompson series is
\[
\Tg(\tau)\;=\;\mathrm{tr}_{\Vnat}\bigl(g\,q^{L_0-1}\bigr),
\qquad q=e^{2\pi i\tau},
\]
with $T_e=J$. Automorphisms of $\Vnat$ act by chiral-algebra
automorphisms, hence lift through $\SpCh_{\Sigma_2,C}$ to
automorphisms of $\PhiFA_3(X_B)$ covering the Leech-orbifold
cover $X_B\to X_B$.

\begin{proposition}[$g$-twisted second stage]
\label{prop:g-twisted-spec}
For each $g\in\Mon$ with lift $\hat g$ to
$\Aut\PhiFA_3(X_B)$, the second-stage specialisation at the
$\hat g$-twisted sector $(\Sigma_2,C)_{\hat g}$ gives the $g$-twisted
module $\Vnat(g)$, whose graded character on $E_\tau$ is
$\Tg(\tau)$. The Zhu functor on the $g$-twisted module recovers
the FLM/Dong--Mason twisted-trace.
\end{proposition}

The lift $g\mapsto\hat g$ uses the Monster-orbifold pattern: an
element of $\Mon$ acts on the Leech part $\Sigma_2$ through its
Conway--Niemeier representative in $\mathrm{Co}_0$, and on the
$E_\tau$-factor by a fractional linear transformation determined
by the Fricke--Atkin--Lehner level of $g$. The two-stage functor
transports this to the $\hat g$-twisted factorisation
homology integrand.

\section*{3. The Hauptmodul conjecture through two-stage factorisation}

Conway--Norton 1979 asserts that each $\Tg$ is the
Hauptmodul for a genus-zero subgroup $\Gamma_g\subset\Sp
L_2(\R)^+$ commensurable with $\mathrm{SL}_2(\Z)$; Borcherds 1992
proves this by identifying $\Tg(p)-\Tg(q)$ with a twisted
denominator of $\mMon$.

\begin{theorem}[Hauptmodul property as two-stage rigidity]
\label{thm:hauptmodul-two-stage}
\ClaimStatusTheorem{} \emph{(Borcherds 1992 Thm.~1.1; restated.)}
For every $g\in\Mon$, $\Tg$ is a Hauptmodul for a genus-zero
$\Gamma_g$. In two-stage terms: the map
\[
\Xi_g\;:\;E_\tau\longmapsto
\SpCh_{(\Sigma_2)_{\hat g},\,E_\tau}\bigl(\PhiFA_3(X_B)\bigr)
\]
descends to an isomorphism $\Xi_g\colon
X_0(\Gamma_g)\xrightarrow{\sim}\mathbb{P}^1$ after passing to
graded characters. Genus zero of $X_0(\Gamma_g)$ is the
\emph{rigidity} statement that the $\hat g$-twisted
$E_3$-hFA on $E_\tau$ is determined by a single modular function
$\Tg$.
\end{theorem}

\begin{proof}[Sketch]
Borcherds 1992 constructs from $\PhiFA_3(X_B)$ the Monster Lie
algebra $\mMon$ with twisted denominator
\(
\Tg(p)-\Tg(q)=p^{-1}\prod_{m>0,n}(1-p^m q^n)^{c_g(mn,n)},
\)
where $c_g(mn,n)$ are Fourier coefficients of $\Tg$. The Euler
character on both sides plus the Borcherds BKM Weyl identity
forces $\Tg$ to be a replicable function of genus zero
(Borcherds 1992 \S\S 7--8). In the two-stage picture the BKM
denominator is the logarithm of the stage-2 partition function at
the $\hat g$-twisted cusp; rigidity follows because
$\SpCh$ at a codim-$1$ cusp of the $\IIonel$ period domain is
parametrised by a single cross-ratio $q$.
\end{proof}

\section*{4. Carnahan 2012 generalised moonshine as
commuting-pair second stage}

Norton's 1987 generalised-moonshine conjecture, proved by
Carnahan 2012 (Thm.~1.4), assigns to each commuting pair
$(g,h)\in\Mon\times\Mon$ a modular function $\Tgh$ with the
properties: (i) $\Tgh$ depends only on the conjugacy class of
$(g,h)$ in $\Mon$ up to scalar; (ii) under
$\bigl(\begin{smallmatrix}a&b\\c&d\end{smallmatrix}\bigr)\in\Sp
L_2(\Z)$,
\(
T_{(g^ah^b,g^ch^d)}=\gamma\cdot\Tgh
\)
for a root of unity~$\gamma$; (iii) each $\Tgh$ is either constant
or a Hauptmodul. Carnahan's proof proceeds through Borcherds-product
lifts of the $(g,h)$-twisted twined characters of $\Vnat$
(arXiv:$0812.3440$ \S\S 5--7).

\begin{proposition}[Two-stage hosting of generalised moonshine]
\label{prop:two-stage-hosting}
Fix a commuting pair $(g,h)\in\Mon\times\Mon$. The orbifold
\(
X_{B,(g,h)}=X_B/\langle\hat g,\hat h\rangle
\)
is a CY$_3$ with an $\langle\hat g,\hat h\rangle$-twisted
$E_3$-hFA $\PhiFA_3(X_{B,(g,h)})$, and
\[
\SpCh_{(\Sigma_2)_{\hat g},\,E_\tau/\hat h}
\bigl(\PhiFA_3(X_{B,(g,h)})\bigr)
\;=\;\Vnat(g,h) \quad\text{as $h$-twisted $g$-twisted module},
\]
whose graded character on $E_\tau/\hat h$ is $\Tgh(\tau)$. The
Carnahan $\Sp L_2(\Z)$-equivariance is the stage-$1$ automorphism
covariance
\(
\PhiFA_3\bigl(\sigma\cdot X_{B,(g,h)}\bigr)
=\sigma^*\PhiFA_3(X_{B,(g,h)})
\)
transported through the second-stage modular-curve specialisation;
the scalar~$\gamma$ is the Dijkgraaf--Witten anomaly cocycle of the
$\langle\hat g,\hat h\rangle$-orbifold of the Leech sigma model.
\end{proposition}

The $(g,h)$-structure is the natural $(0+1)$-dimensional echo of
the $E_3$-hFA: the commuting pair labels a point of
$\mathrm{Hom}(\pi_1(E_\tau),\Mon)/\Mon$, so $\Tgh$ is the
partition function of $\PhiFA_3(X_B)$ on the flat
$\Mon$-bundle over $E_\tau$ with holonomies $(g,h)$. Carnahan's
Hauptmodul conclusion is stage-1 rigidity of
$\PhiFA_3(X_B)$ transported to stage-2 by the same cusp analysis
as Theorem~\ref{thm:hauptmodul-two-stage}. The Duncan 2007
super-moonshine enhancement (Conway-moonshine $V^{f\natural}$) sits
in the same picture with $\Leech$ replaced by the odd Leech
lattice and $\Sigma_2$ replaced by its super-analogue.

\section*{5. Consequence: $j(\tau)$ as universal-CY$_3$ witness}

Assemble the data. The $E_3$-hFA $\PhiFA_3(X_B)$ on the
single CY$_3$ $X_B$ produces, through the second-stage functor
$\SpCh_{\Sigma_2,C}$ applied to the Leech-$\mathrm{Mon}$-orbifold
tower:
\begin{itemize}
\item the Monster VOA $\Vnat$ at $(\Sigma_2,C)=(T^{24}_\Leech/\Z_2,
E_\tau)$, character $J(\tau)=j(\tau)-744$;
\item each McKay--Thompson $\Tg$ at the $\hat g$-twisted cusp;
\item each Carnahan $\Tgh$ at the commuting-pair orbifold;
\item the Monster BKM denominator $j(p)-j(q)$ at the
$(T^{24}_\Leech,\mathrm{pt})$ cusp (note S3);
\item the Fake-Monster BKM at $(T^{24}_\Leech,E_\tau)$ (note S4).
\end{itemize}

\begin{conjecture}[$j(\tau)$ is a two-stage witness of CY$_3$ structure]
\label{conj:j-universal-witness}
The Klein $j$-function is the
$q^0$-partition function of the trivial specialisation of
$\PhiFA_3(X_B)$:
\[
j(\tau)-744\;=\;Z_{E_\tau}\bigl(\SpCh_{T^{24}_\Leech/\Z_2,\,E_\tau}
\PhiFA_3(X_B)\bigr).
\]
The entire moonshine module $\Vnat$, the McKay--Thompson family
$\{\Tg\}_{g\in\Mon}$, the Carnahan family $\{\Tgh\}$, and the
Borcherds denominators $\{\Tg(p)-\Tg(q)\}$ are simultaneous
specialisation witnesses of the single canonical $E_3$-hFA
$\PhiFA_3(X_B)$. The $\Sp L_2(\Z)$-rigidity of each witness is
controlled by the codim-$1$ cusp structure of the period domain
$\mathcal{D}_{\IIonel}$; the Hauptmodul property is the stage-$1$
uniqueness of $\PhiFA_3$ up to contractible choice.
\end{conjecture}

This is the Vol III universal reading: moonshine is not a coincidence
among three genus-zero lists, but the shadow on $E_\tau$ of a single
$E_3$-hFA on $X_B$, evaluated at the Leech-$\Mon$-orbifold
tower of specialisation data. The arithmetic scalars
$\kBKM(\mathrm{Monster})=0$, $\kBKM(\Delta_5)=5$, and the
spectrum $\{2,3,5,24\}$ on $K3\times E$ are second-stage cusp
multiplicities; the Hauptmodul property of every $\Tg$ is the
first-stage rigidity of $\PhiFA_3$.

\paragraph{Primary sources.} Borcherds 1986 \emph{Proc.\ Natl.\ Acad.\
Sci.}\ \textbf{83} (vertex algebras); Borcherds 1992 \emph{Invent.\ Math.}\
\textbf{109} (monstrous moonshine, Thms.~1.1, 6.1, \S 8 twisted
denominator); Frenkel--Lepowsky--Meurman 1988 \emph{VOA and the Monster}
\S\S 10--12 (Leech $\Z/2$-orbifold, Thm.~12.3.1); Conway--Norton 1979
\emph{Bull.\ LMS}\ \textbf{11} (Hauptmodul conjecture); Carnahan 2012
(arXiv:$0812.3440$, Thm.~1.4, generalised moonshine, Borcherds-product
construction of $\Tgh$); Duncan 2007 \emph{Math.\ Res.\ Lett.}\
\textbf{14} (super-moonshine $V^{f\natural}$); Costello--Gwilliam 2017
(holomorphic factorisation algebras, $E_3$-hFA on curves).

\end{document}
